\section{Critical recovery with growing rank}
\label{sec:critical-rank}
Consider two disjoint column bases on a quotient of rank $r$. Throughout this section $\|\pi\|_{15/7}\le B$ and logarithms are natural. We give two sufficient conditions at the critical moment. One uses only the quotient rank and the matching order. The other works for larger ranks if the exceptional syndromes satisfy sufficiently many linear restrictions.

\paragraph{Theorem.}
Let $k\to\infty$ and impose the matching hypotheses of (F1). For arbitrary two-base quotients with $r=o(k^{3/2})$ the entire ordered spectrum recovers that of ordinary noise, including multiplicities. More precisely, if $A_t=\{U>e^t\}$ and $s_t$ is the norm of the chain killed outside $A_t$, then for all $t>0$
\[
 s_t\le \exp\{-c t^2/(r+t)+C_B(1+r/t)\}.                  \tag{G1}
\]
This bound for the killed chain requires no matching. For the second condition, suppose the exceptional syndromes in (F1) lie in the inverse image of at most $L$ values of a rank-$d$ linear map on the syndrome space. Put $\kappa=(2/5)\log 2-\log(5/4)>0$. Its killed norm satisfies
\[
 s\le \exp\{-\kappa d+(2/5)\log L+
             C_B[1+(r-d)^{1/3}+\sqrt d]\}.                  \tag{G2}
\]
In particular a single affine subspace with $d/r^{1/3}\to\infty$ suffices for ordered recovery, even when $d/r\to0$. More generally one can take $L\le 2^{\beta d}$ with fixed $\beta<1-(5/2)\log_2(5/4)$. The constants and powers give sufficient regimes. Sharpness of these regimes is open.

\paragraph{Reduction to two ordinary product operators.}
On a union of syndrome fibers the top killed eigenvalue agrees with that of the quotient. Indeed the symmetric compression $M_{\sqrt{U}1_E} P M_{\sqrt{U}1_E}$ has a nonnegative top eigenvector and commutes with translations in the kernel of the quotient map. Averaging that vector gives a nonzero fiber-constant top eigenvector when the norm is positive. A further restriction inside fibers only decreases the top eigenvalue by the variational principle. The same argument applies to the unrestricted weighted compression for every nonnegative quotient weight $W$. Write the quotient noise as $P_1P_2J$, where $P_i$ is product noise of correlation $1/2$ in column basis $i$, and $J$ contains unused columns. The three commuting convolution operators have Fourier multipliers in $[0,1]$. Thus for any vector $z$,
\[
 \langle z,P_1P_2Jz\rangle
 \le\langle z,P_1P_2 z\rangle\le\|P_1z\|_2\|P_2z\|_2 .
\]
All quotient function norms use uniform probability. Taking $z=\sqrt W f$, duality and H\"older show the weighted compression has norm at most
\[
 \|W\|_v\prod_{i=1}^2\|1_{\{W>0\}}P_i\|_{2\to m},
 \qquad m=2/(1-1/v),\quad v>1.                             \tag{G3}
\]

\paragraph{The cost of increasing the one-bit output order.}
For one-bit half noise $N$, $\|N\|_{2\to5+\epsilon}\le
\exp(C\epsilon^2)$ for $0\le\epsilon\le\epsilon_0$. By positivity, scaling and reflection take input $(1+x,1-x)$ for $0\le x\le1$. At output order five the fifth power of the output norm is $1+5x^2/2+5x^4/16$. The fifth power of the input norm is $(1+x^2)^{5/2}\ge1+5x^2/2+15x^4/8$. Since both expressions are uniformly bounded, the log norm ratio is at most $-c_0 x^4$. For $m$ in a compact neighborhood of five,
\[
 \partial_m\big[m^{-1}\log(((1+x/2)^m+(1-x/2)^m)/2)\big]=O(x^2).
\]
The expression on the left is even, vanishes at zero and has bounded second $x$ derivative on the indicated compact rectangle. Maximizing $-c_0x^4+C\epsilon x^2$ proves the estimate. Minkowski's mixed-norm inequality, first in two factors and then by induction, bounds the norm on $j$ bits by $\exp(Cj\epsilon^2)$ as $m\ge2$. This calculation refines the usual two-point Bonami bound immediately beyond its equality exponent. Related concentration inequalities for unweighted cube noise are due to \href{https://arxiv.org/abs/1606.07491}{Polyanskiy and Samorodnitsky}.

\paragraph{The high-multiplier estimate.}
For $0<\delta\le\delta_0$ small, let $q=15-\delta$, $a=7/15+2/q$, and $v_0=1/a$. The potential estimate in the fractional-packing section gives explicitly
\[
 \|g^{-1}\|_q\le 8B(16B^2/3)^{16/\delta},
 \qquad \|U\|_{v_0}\le \|\pi\|_{15/7}\|g^{-1}\|_q^2
                         \le \exp(C_B/\delta).
\]
Indeed $h\ge1/(8B)$ on measure at least $3/(16B^2)$, and reverse hypercontractivity starts from order $\delta/16$. To exploit the cutoff take $1/v=a+\delta$. On $A_t$,
\[
 \|U1_{A_t}\|_v\le e^{-(v_0/v-1)t}\|U\|_{v_0}^{v_0/v}
               \le \exp(-c_1\delta t+C_B/\delta).
\]
Here $v_0/v-1=\delta/a$, and $5\le m\le 5+C\delta$. Dropping the support restrictions in (G3) now gives a log norm bound of $-c_1\delta t+C_B/\delta+Cr\delta^2$. Set $\delta=\eta t/(r+t)$ with fixed sufficiently small $\eta>0$. This proves (G1). Put $t=\sigma k$ from (F1). The negative term dominates when $r=o(k^{3/2})$, and $t^2/(r+t)$ tends to infinity. The upper half of (F1) yields the assertion.

\paragraph{Affine lists, including diffuse syndromes.}
In coordinates of either basis choose $d$ pivot coordinates for the given linear map. For each assignment of the remaining coordinates there are at most $L$ acceptable pivot strings. With $E$ the whole list, the product noise in the pivot coordinates maps $L^2$ to $L^\infty$ with norm $(5/4)^{d/2}$: the one-bit row density relative to uniform has squared norm $5/4$. In the other coordinates commute the noise past the pivot noise. Positivity and the triangle inequality bound the pivot $L^2$ norm after complementary noise by that noise applied to the input's pivot $L^2$ norm. Integrating on the allowed strings thus gives in either basis
\[
 \|1_EP_i\|_{2\to m}
 \le (L2^{-d})^{1/m}(5/4)^{d/2}\exp(C(r-d)(m-5)^2).
\]
Now take $v=v_0$ in (G3), so $2/m=2/5-O(\delta)$, and use $\|U1_E\|_{v_0}\le \exp(C_B/\delta)$. The logarithm of the resulting bound is at most
\[
 -\kappa d+(2/5)\log L+C_B/\delta+C d\delta+C(r-d)\delta^2 .
\]
Choose $\delta=\delta_0\min((1+r-d)^{-1/3},(1+d)^{-1/2})$. This proves (G2). With $d=r$ the estimate recovers the syndrome-count threshold. For $L=1$ and $r^{1/3}\ll d\ll r$, the affine subspace has $2^{r-d}$ elements, exponentially more than that count threshold permits.

\paragraph{Conditional mass and scope.}
For any stationary cut $F$ take the counting-space unit vector $f=1_F\sqrt{\nu/\nu(F)}$, where $\nu=\pi\mu$. With $T=M_{\sqrt U}P M_{\sqrt U}$ the return probability is $R=\langle f,Tf\rangle$. For the multiplication projection $D$ onto an affine list, positivity and $\|T\|\le1$ give $\sqrt R\le\sqrt{\|DTD\|}+\|(1-D)f\|$. Consequently a return bounded below by $s_0>0$ forces at least $s_0-o(1)$ of the conditional mass outside every list for which (G2) tends to zero. Taking $d=r$, the largest $\lfloor 2^{\beta r}\rfloor$ atoms give the Shannon entropy bound $\liminf H_2(A X\mid X\in F)/r\ge s_0[1-(5/2)\log_2(5/4)]$ as $r\to\infty$, where $H_2$ here denotes Shannon entropy in bits and $A$ the quotient map. Each atom outside the list has mass at most $1/\lfloor 2^{\beta r}\rfloor$; let $\beta$ increase to the threshold. Unrestricted critical recovery at arbitrarily large rank remains open.
